\exercisesfor{7}

The conventions of § 7 remain valid unless otherwise mentioned.

\begin{enumerate}
\def\labelenumi{\arabic{enumi}.}
\tightlist
\item
  \begin{enumerate}
  \def\labelenumii{(\alph{enumii})}
  \tightlist
  \item
    For a Lie K-algebra to be nilpotent, it is necessary and sufficient that it be isomorphic to a subalgebra of an algebra \(n(n, \mathbf{K})\) .
  \end{enumerate}
\end{enumerate}

\begin{enumerate}
\def\labelenumi{(\alph{enumi})}
\setcounter{enumi}{1}
\tightlist
\item
  Suppose that K is algebraically closed. For a Lie K-algebra to be solvable, it is necessary and sufficient that it be isomorphic to a subalgebra of an algebra \(t(n, \mathbf{K})\) .
\end{enumerate}

\begin{enumerate}
\def\labelenumi{\arabic{enumi}.}
\setcounter{enumi}{1}
\item
  Let g be a Lie algebra and n its largest nilpotent ideal. There exist a finite-dimensional vector space V and an isomorphism of g onto a subalgebra of \(\mathfrak{sl}(V)\) , which maps every element of n to a nilpotent endomorphism of V. (Use Ado's Theorem and Exercise 5 of § 1.)
\item
  Let g be a Lie algebra and U its enveloping algebra.
\end{enumerate}

\begin{enumerate}
\def\labelenumi{(\alph{enumi})}
\item
  Show that for all \(u \in \mathrm{U}\) ( \(u \neq 0\) ) there exists a finite-dimensional representation \(\pi\) of \(\mathrm{U}\) such that \(\pi(u) \neq 0\) . (Use Exercise 2 and Exercise 5 (b) of § 3.)
\item
  Deduce from (a) that if g is semi-simple there exists an infinity of inequivalent finite-dimensional simple representations of g. (Let \(\rho_{1},\ldots,\rho_{n}\) be finite-dimensional simple representations of g, \(N_{1},\ldots,N_{n}\) the annihilators of the corresponding U-modules and \(N=\bigcap_{i=1}^{n}N_{i}\) . Show that N is of finite co-dimension in U. Let \(n\in N,n\neq0\) . Apply (a) to n.)
\end{enumerate}

\begin{enumerate}
\def\labelenumi{\arabic{enumi}.}
\setcounter{enumi}{3}
\tightlist
\item
  Let g be a Lie algebra, a a subinvariant subalgebra of g, \(\mathfrak{r}(a)\) the radical of a and \(\rho\) a finite-dimensional representation of a. For \(\rho\) to admit a finite-dimensional extension to g, it is necessary and sufficient that \(\mathcal{D}_{\mathfrak{g}} \cap \mathfrak{r}(a)\) be contained in the largest nilpotency ideal of \(\rho\) . (For the sufficiency proceed as follows; let \(g = g_{0} \supset g_{1} \supset \cdots \supset g_{k} = a\) be a decomposition series of g with the properties of Exercise 11 (b) of §6. Show by induction the existence of an extension \(\rho_{i}\) of \(\rho\) to \(g_{0}\) , whose largest nilpotency ideal contains \(\mathcal{D}_{\mathfrak{g}} \cap \mathfrak{r}(g_{i})\) . In the notation of Exercise 11 (b) of §6, the passage from \(\rho_{i+1}\) to \(\rho_{i}\) follows from Theorem 1 when \(b_{i}\) is simple, or when \(b_{i}\) is 1-dimensional and
\end{enumerate}

\[
\mathscr {D} \mathfrak {g} \cap \mathfrak {r} (\mathfrak {g} _ {i + 1}) = \mathscr {D} \mathfrak {g} \cap \mathfrak {r} (\mathfrak {g} _ {i}).
\]

When \(\mathcal{D}\mathfrak{g} \cap \mathfrak{r}(\mathfrak{g}_{i+1}) \neq \mathcal{D}\mathfrak{g} \cap \mathfrak{r}(\mathfrak{g}_i)\) , it can be assumed that \(\mathfrak{h}_i \subset \mathcal{D}\mathfrak{g} \cap \mathfrak{r}(\mathfrak{g}_i)\) . Then \(\mathfrak{r}(\mathfrak{g}_i) = \mathfrak{r}(\mathfrak{g}) \cap \mathfrak{g}_i\) (§ 5, Exercise 10) and Theorem 1 can be applied.)

¶ 5. Let K be a field of characteristic p \textgreater{} 0, g a Lie algebra of finite dimension n over K and U the enveloping algebra of g.

\begin{enumerate}
\def\labelenumi{(\alph{enumi})}
\item
  A \(p\) -polynomial (in one indeterminate) over \(\mathbf{K}\) is any polynomial in \(\mathrm{K}[\mathrm{X}]\) whose only terms \(\neq 0\) are of degree a power of \(p\) . Show that for every polynomial \(f(\mathbf{X}) \neq 0\) in \(\mathrm{K}[\mathrm{X}]\) there exists \(g(\mathbf{X}) \in \mathrm{K}[\mathbf{X}]\) such that \(f(\mathbf{X})g(\mathbf{X})\) is a \(p\) -polynomial (consider the remainders in Euclidean division of the \(\mathbf{X}^{p_k}\) by \(f(\mathbf{X})\) ).
\item
  Show that for every element \(z \in \mathfrak{g}\) there exists a non-zero polynomial \(f(\mathbf{X}) \in \mathbf{K}[\mathbf{X}]\) such that \(f(z)\) belongs to the centre C of U. (Consider the minimal polynomial of the endomorphism ad \(z\) and apply (a) to this polynomial, and also formula (1) of Exercise 19 of § 1.)
\item
  Let \((e_i)_{1\leqslant i\leqslant n}\) be a basis of \(\mathfrak{g}\) . For each \(i\) let \(f_{i}\neq 0\) be a polynomial in \(\mathrm{K}[\mathrm{X}]\) such that \(f_{i}(e_{i})\in \mathrm{C}\) ; let \(d_{i}\) be its degree and \(\mathrm{L}\) the two-sided ideal of \(\mathrm{U}\) generated by the \(y_{i} = f_{i}(e_{i})\) . Show that the classes mod. \(\mathrm{L}\) of the elements \(e_1^{\alpha_1}\ldots e_n^{\alpha_n}\) , where \(0\leqslant \alpha_{i} < d_{i}\) form a basis of \(\mathrm{U / L}\) .
\item
  Show that \(\mathfrak{g}\) admits a finite-dimensional faithful linear representation (choose the \(f_{i}\) so that \(d_{i} > 1\) for all \(i\) ).
\item
  Show that if \(g \neq \{0\}\) there exists a non-semi-simple finite-dimensional linear representation of g. (Assuming that the \(f_{t}\) are of degrees \(d_{t} > 1\) , replace the \(f_{t}\) by \(f_{t}^{2}\) in the definition of L and note that there then exist nilpotent elements \(\neq 0\) in the centre of U/L and hence that U/L is not semi-simple.)
\end{enumerate}


